\subsection*{Commutativity}
\label{subsec:commutativity}

\begin{tcolorbox}[title={Definition --- Commutativity}]
\begin{definition}[Commutativity]
\label{def:commutativity}
Let $\star$ be a binary operation on $G$.
The operation $\star$ is \textbf{commutative} if
\[
a \star b = b \star a
\quad \text{for all } a,b \in G.
\]
\end{definition}
\end{tcolorbox}

\begin{remark*}[Standard quantified statement]
\[
\operatorname{Commutative}(\star,G)
\;\colonequiv\;
\forall a,b \in G,\; a \star b = b \star a.
\]
\end{remark*}

\begin{remark*}[Definition predicate reading]
$\star$ is commutative on $G$ exactly when swapping the two inputs never changes
the output.
\end{remark*}

\begin{remark*}[Negated quantified statement]
\[
\neg\operatorname{Commutative}(\star,G)
\;\colonequiv\;
\exists a,b\in G,\; a\star b\neq b\star a.
\]
\end{remark*}

\begin{remark*}[Interpretation]
Commutativity says that the order of the two inputs does not affect the
result.  Associativity and commutativity are independent properties.
\end{remark*}

\begin{dependencies}
\begin{itemize}
  \item Binary operation.
\end{itemize}
\end{dependencies}
